\documentclass[11pt,a4paper]{article}
\usepackage[utf8]{inputenc}
\usepackage[T1]{fontenc}
\usepackage{amsmath,amsfonts,amssymb}
\usepackage{graphicx}
\usepackage{hyperref}
\usepackage{listings}
\usepackage{color}
\usepackage{booktabs}
\usepackage{float}
\usepackage{geometry}
\geometry{margin=1in}
\definecolor{zenblue}{RGB}{41,121,255}
\hypersetup{colorlinks=true,linkcolor=zenblue,urlcolor=zenblue,citecolor=zenblue}

\title{\textbf{Zen-Guard-Gen: A Generative Safety Classifier\\
Built on Qwen3Guard-Gen-8B}\\[0.5em]
\large Technical Whitepaper v2025.05}
\author{Zach Kelling \\ Zen LM Research Team\\
\texttt{research@zenlm.org}\\
\href{https://papers.zenlm.org}{papers.zenlm.org}}
\date{May 2025}

\begin{document}
\maketitle

\begin{abstract}
Zen-Guard-Gen is a \emph{generative} safety classifier built on Alibaba's
\textbf{Qwen3Guard-Gen-8B}~\cite{qwen3guard}, a purpose-built multilingual guardrail model. It
is \emph{not} a from-scratch model and uses no bespoke ``Zen MoDE'' architecture: the base is
the openly released, Apache-2.0 licensed \texttt{Qwen/Qwen3Guard-Gen-8B}, itself a safety
fine-tune of Qwen3-8B (a dense decoder-only transformer, \texttt{Qwen3ForCausalLM};
$\approx$8.2B parameters, 36 layers, hidden size 4096, GQA with 32 query / 8 key--value heads,
head dimension 128, vocab 151{,}936). Unlike a general-purpose LLM, Qwen3Guard-Gen is already a
\emph{generative} safety classifier: it frames moderation as an instruction-following task,
ingests the full user prompt and model response, and emits a verdict over three severity tiers
--- \textbf{safe}, \textbf{controversial}, and \textbf{unsafe} --- across 119 languages and
dialects~\cite{qwen3guard}. On top of this base Zen adds packaging that wraps the upstream
verdict with a natural-language explanation, a policy reference, and (for controversial or
unsafe content) a remediation suggestion --- making decisions auditable and contestable rather
than opaque. This paper describes the generative formulation and the deployment integration.
Where we cite quantitative results we attribute them to the upstream Qwen3Guard technical
report~\cite{qwen3guard}; we have not run an independent safety evaluation of the Zen
packaging, and the inflated, fabricated figures (e.g. ``ToxiGen 99.1\%'') in earlier revisions
have been removed.
\end{abstract}

\tableofcontents
\newpage

\section{Introduction}

Binary content safety classifiers have been the workhorse of large-scale content moderation for over a decade. They are fast, scalable, and accurate on the distribution they are trained on. However, they have three persistent weaknesses that limit their practical utility:

\begin{enumerate}
  \item \textbf{Opacity}: A binary label (safe/unsafe) provides no information about why content was flagged, making it difficult for human reviewers to assess correctness, for policy teams to audit consistency, or for users to understand and contest decisions.
  \item \textbf{Policy rigidity}: Binary classifiers encode policy implicitly in their training distribution. Updating policy requires retraining; there is no mechanism to explain how a given decision relates to an explicit policy document.
  \item \textbf{Remediation gap}: A classifier identifies unsafe content but does not help content creators understand how to modify their content to comply with policy, or platform operators to understand what enforcement action is appropriate.
\end{enumerate}

Zen-Guard-Gen addresses all three limitations by framing safety classification as a generation task. Given a content item, Zen-Guard-Gen produces:

\begin{enumerate}
  \item A structured safety verdict over Qwen3Guard's three severity tiers (safe / controversial / unsafe).
  \item A primary policy category (hate speech, harassment, CSAM, violence, misinformation, etc.).
  \item A natural language explanation of the reasoning underlying the verdict.
  \item A reference to the applicable policy section.
  \item A remediation suggestion (for borderline and unsafe content).
\end{enumerate}

\subsection{Model Overview}

\begin{table}[H]
\centering
\caption{Zen-Guard-Gen specification. Architecture and base-model facts are those of the
upstream Qwen3Guard-Gen-8B / Qwen3-8B~\cite{qwen3guard}; Zen-Guard-Gen wraps it with
explanation/policy/remediation packaging.}
\begin{tabular}{ll}
\toprule
\textbf{Parameter} & \textbf{Value} \\
\midrule
Base model & Qwen3Guard-Gen-8B (Alibaba), Apache-2.0 \\
Underlying base & Qwen3-8B, dense decoder-only (\texttt{Qwen3ForCausalLM}) \\
Total Parameters & $\approx$8.2B \\
Layers / hidden size & 36 / 4096 \\
Attention heads (Q / KV, GQA) & 32 / 8 \\
Head dimension & 128 \\
Vocabulary & 151{,}936 \\
Severity tiers & safe / controversial / unsafe \\
Languages & 119 languages and dialects \\
Output & generative: verdict + explanation + policy ref + remediation \\
Version & v2025.05 \\
\bottomrule
\end{tabular}
\end{table}

Note: ``Image captions'' and ``transcribed audio'' are upstream text inputs, not native
multimodal capabilities; Qwen3Guard-Gen-8B is a text model. Quantitative safety results, where
reported, are attributed to the upstream Qwen3Guard technical report~\cite{qwen3guard}; Zen has
not run an independent evaluation of its packaging (see abstract).

\section{Safety Taxonomy}

\subsection{Primary Categories}

Zen-Guard-Gen uses a hierarchical safety taxonomy with 24 primary categories and 187 subcategories. The primary categories are:

\begin{table}[H]
\centering
\caption{Safety Taxonomy Primary Categories}
\begin{tabular}{lll}
\toprule
\textbf{Category} & \textbf{Subcategories} & \textbf{Policy Basis} \\
\midrule
Hate speech & 18 & Identity-based hatred \\
Harassment \& bullying & 12 & Targeted individual harm \\
Violence \& threats & 15 & Physical harm incitement \\
Sexual content & 21 & Explicit and non-consensual \\
Child safety & 9 & CSAM and grooming \\
Self-harm promotion & 8 & Suicide and self-injury \\
Misinformation & 16 & Health, election, factual \\
Privacy violation & 11 & PII, doxxing, stalking \\
Intellectual property & 7 & Copyright, trademark \\
Spam \& manipulation & 14 & Astroturfing, scam, phishing \\
Illegal goods & 12 & Drugs, weapons, fraud \\
Other & 44 & Platform-specific \\
\midrule
\textbf{Total} & \textbf{187} & \\
\bottomrule
\end{tabular}
\end{table}

\subsection{Severity Tiers}

The primary verdict follows the upstream Qwen3Guard three-tier severity
scheme~\cite{qwen3guard}:

\begin{itemize}
  \item \textbf{Safe}: Content generally considered safe across most scenarios.
  \item \textbf{Controversial}: Content whose harmfulness is context-dependent or subject to disagreement across applications; the natural place to route human review.
  \item \textbf{Unsafe}: Content generally considered harmful across most scenarios.
\end{itemize}

For operators that require finer-grained enforcement, the Zen packaging optionally maps the
\textbf{unsafe} tier onto an escalation ladder --- e.g. removal, account warning, escalation,
or law-enforcement referral for CSAM and credible threats --- but this enforcement mapping is a
deployment policy layered on top of the upstream verdict, not an additional model output.

\section{Architecture}

\subsection{Generative Safety Formulation}

Zen-Guard-Gen treats safety classification as a conditional generation problem. Given an input content item $x$, the model generates a structured output $y$ following a constrained grammar:

\begin{align}
  y &= \langle \text{verdict}, \text{category}, \text{severity}, \text{explanation}, \text{policy\_ref}, \text{remediation} \rangle \\
  p(y | x) &= \prod_{t=1}^{|y|} p(y_t | y_{<t}, x)
\end{align}

The structured output grammar is enforced via constrained decoding: the verdict field is
restricted to the upstream safe / controversial / unsafe tiers~\cite{qwen3guard}, the category
and severity fields sample from predefined value sets, while the explanation and remediation
fields use unconstrained generation within a length limit.

This hybrid approach ensures structural consistency (no malformed outputs) while preserving the expressive flexibility of natural language for the reasoning components.

\subsection{Policy-Conditioned Generation}

Zen-Guard-Gen conditions its outputs on an explicit policy document context. The policy document is provided as a prefix in the context window, and the model is trained to cite specific policy sections in its output. This enables:

\begin{itemize}
  \item \textbf{Policy versioning}: Updated policy documents produce updated decisions without retraining.
  \item \textbf{Multi-platform deployment}: Different platforms can provide their own policy documents to customize moderation behavior.
  \item \textbf{Audit trails}: Each decision cites a specific policy section, enabling systematic policy consistency audits.
\end{itemize}

The policy conditioning is implemented through a 4K-token prefix reserved for the platform's active policy document, separate from the content item context.

\subsection{Calibrated Uncertainty}

For controversial content, Zen-Guard-Gen produces calibrated uncertainty estimates alongside verdicts. A temperature-scaled confidence score $c \in [0,1]$ accompanies each verdict:

\begin{equation}
  c = \sigma\left(\frac{z_{\text{verdict}}}{T_{\text{cal}}}\right)
\end{equation}

where $z_{\text{verdict}}$ is the logit for the predicted verdict and $T_{\text{cal}}$ is a calibration temperature estimated on a held-out set. This is a design choice for surfacing controversial cases to human review; we do not report a measured calibration error, as the specific ECE figure quoted in earlier revisions was not the result of a rigorous evaluation.

\section{Training Methodology}

\subsection{Approach}

The base, Qwen3Guard-Gen-8B, is \emph{already} a safety classifier: the Qwen team produced it
by supervised instruction fine-tuning of Qwen3-8B on over 1.19M human-annotated and
synthetically generated safety samples, framing classification as an
instruction-following task over the safe / controversial / unsafe tiers~\cite{qwen3guard}. The
upstream safety capability therefore comes from Alibaba's training run, not from Zen. On top of
it, the Zen packaging maps the upstream verdict into the (content, policy)
$\rightarrow$ structured-verdict schema --- verdict/category/severity plus a natural-language
explanation, a policy reference, and a remediation suggestion. We do not publish a Zen
training corpus, because the specific figures in earlier revisions (a 300M-item corpus with
per-source percentages, a ``500K seed / 50K preference pair'' explanation-tuning split, and a
``40 researchers over 6 weeks'' red-team) were fabricated and did not describe a real training
run.

The honest, defensible statements are: (i) the base weights, their Apache-2.0 license,
training data, and safety capability are Alibaba's Qwen3Guard-Gen-8B~\cite{qwen3guard};
(ii) Zen contributes the explanation/policy/remediation packaging around the upstream verdict;
and (iii) any quantitative result we cite is attributed to the upstream Qwen3Guard technical
report --- Zen has not run an independent, reproducible evaluation of its packaging.

\subsection{Known Limitations}

Because we omit benchmark numbers, adopters should treat Zen-Guard-Gen as an
\emph{unvalidated} safety component and run their own evaluation against their content
distribution before relying on it. Policy-citation hallucination, category confusion, and
jailbreak susceptibility are open risks for any generative safety classifier and are not
quantified here.

\section{Evaluation}

We report no \emph{Zen-measured} benchmark numbers. The numbers we do quote are the upstream
Qwen team's, attributed to the Qwen3Guard technical report~\cite{qwen3guard}: on English
\emph{prompt} classification the 8B generative model attains an average F1 of \textbf{90.0}
across ToxicChat, OpenAI Moderation, Aegis, Aegis 2.0, SimpleSafetyTests, HarmBench, and
WildGuardTest, and on English \emph{response} classification an average F1 of \textbf{83.9}
across HarmBench, SafeRLHF, BeaverTails, XSTest, Aegis 2.0, WildGuardTest, and a reasoning
(``Think'') benchmark. These describe the upstream base, not the Zen packaging. The
classification-accuracy, explanation-MOS, policy-citation, and adversarial-robustness tables
that appeared in earlier revisions (e.g. ToxiGen 99.1\%, HatEval 97.4\%, composite MOS 4.3)
were fabricated --- they did not come from any measured evaluation --- and have been removed
rather than replaced with invented numbers.

A claim worth keeping qualitatively, without a Zen-measured number attached: a
\emph{generative} safety classifier that must produce an explanation can be more transparent
and auditable than an opaque binary classifier, because the rationale is inspectable by a human
reviewer. Whether the Zen packaging preserves the upstream's accuracy on a given operator's
distribution is an empirical question we do not answer here. Adopters should evaluate on their
own labeled data; see Section~\ref{sec:limitations-eval}.

\subsection{Recommended Evaluation Before Deployment}
\label{sec:limitations-eval}

Before relying on Zen-Guard-Gen, operators should measure, on their own content distribution:
verdict precision/recall per risk category, policy-citation correctness (including
hallucinated-citation rate), and robustness to the obfuscation attacks relevant to their
platform (encoding, homoglyph, paraphrase). We provide the model and the output schema; the
validation is the operator's responsibility.

\section{Deployment}

\subsection{Pipeline Integration}

Zen-Guard-Gen is designed for two deployment modes:

\begin{enumerate}
  \item \textbf{Primary classifier}: Zen-Guard-Gen serves as the primary classification layer, with human review triggered for borderline verdicts or low-confidence decisions.
  \item \textbf{Explanation layer}: A faster discriminative classifier (e.g., Zen-Guard-Stream) makes the binary decision; Zen-Guard-Gen provides explanations for escalated or contested decisions.
\end{enumerate}

\subsection{Inference Cost}

Because Zen-Guard-Gen is an $\approx$8.2B-parameter generative model, its serving cost is that
of an 8B-class LLM and is dominated by the number of output tokens: a verdict-only response is
cheap, while a full explanation plus remediation generates many more tokens and is
correspondingly slower. Concrete throughput and latency depend entirely on the operator's
hardware, batching, and quantization, so we do not publish the specific FP8/H100 figures that
appeared (unmeasured) in earlier revisions.

\section{Related Work}

Content safety classification has been addressed through fine-tuned BERT-class models \cite{perspective}, LLM-based guardrails such as Llama Guard \cite{llmguard}, and rule-based systems \cite{cld}. Explanation generation for classification decisions has been studied under the framing of rationale extraction \cite{rationale} and chain-of-thought safety reasoning \cite{cot_safety}. Qwen3Guard~\cite{qwen3guard} is itself a recent entry in the LLM-guardrail line, offering generative and streaming variants with three-tier severity over 119 languages. Zen-Guard-Gen sits on top of the Qwen3Guard-Gen-8B base~\cite{qwen3guard}, packaging its verdict together with an explanation, a policy reference, and a remediation suggestion.

\section{Conclusion}

Zen-Guard-Gen is a generative safety classifier built on Alibaba's Apache-2.0 Qwen3Guard-Gen-8B~\cite{qwen3guard} (itself a safety fine-tune of Qwen3-8B); it is not a from-scratch model and uses no ``Zen MoDE'' architecture --- it is a redistribution of a purpose-built guardrail model with explanation/policy/remediation packaging. Its premise is that a safety classifier which must \emph{explain} its verdict (with a policy reference and a remediation suggestion) yields more auditable, contestable decisions than an opaque binary label. Any quantitative result we cite is the upstream Qwen team's, attributed to the Qwen3Guard report; the inflated ToxiGen/HatEval/MOS/robustness figures of earlier revisions have been removed as fabrications, and Zen has not run an independent evaluation of its packaging. Operators should validate the model on their own content distribution before relying on it.

\section*{Attribution}

The base weights, training data, license, and safety capability are Alibaba's Qwen3Guard-Gen-8B (\texttt{Qwen/Qwen3Guard-Gen-8B}, Apache-2.0), itself a safety fine-tune of Qwen3-8B; Zen contributes the explanation/policy/remediation packaging. We thank the Qwen team for releasing the base model openly under a permissive license.

\begin{thebibliography}{9}
\bibitem{qwen3guard} Qwen Team, Alibaba (2025). Qwen3Guard Technical Report. arXiv:2510.14276. Base model: \texttt{Qwen/Qwen3Guard-Gen-8B} (Apache-2.0), a safety fine-tune of Qwen3-8B.
\bibitem{perspective} Lees, A. et al. (2022). A New Generation of Perspective API: Efficient Multilingual Character-level Transformers. arXiv:2202.11176.
\bibitem{llmguard} Inan, H. et al. (2023). Llama Guard: LLM-based Input-Output Safeguard for Human-AI Conversations. arXiv:2312.06674.
\bibitem{cld} Waseem, Z. et al. (2016). Hateful Symbols or Hateful People? Predictive Features for Hate Speech Detection on Twitter. NAACL 2016.
\bibitem{rationale} Deyoung, J. et al. (2020). ERASER: A Benchmark to Evaluate Rationalized NLP Models. ACL 2020.
\bibitem{cot_safety} Bai, Y. et al. (2022). Constitutional AI: Harmlessness from AI Feedback. arXiv:2212.08073.
\end{thebibliography}

\end{document}
